
%%% fhc mu-like

\begin{table}[htb]
\centering
\begin{tabular}{| c | c | c |}
\hline
  {\bf Source of uncertainty}    & \multicolumn{1}{c|}{{\bf ${\delta}N_{SK}$/$N_{SK}$}}\\
                                 & \multicolumn{1}{c|}{(Asimov A)}\\

\hline
\hline
SKDet+FSI+SI                   & 4.13\%                        \\
SKDet only                   & 3.86\%                        \\
FSI+SI only                   & 1.48\%                        \\
\hline
\hline
  Flux                           & 3.60\%                        \\
\hline
\hline
  2p-2h (corr)                     & 3.46\%                        \\
  2p-2h-bar (corr)                  & 0.20\%                        \\
  NC other (uncorr)              & 0.78\%                        \\
  NC 1gamma (uncorr)             & 0.00\%                        \\
  XSec nue/numu (uncorr)                 & 0.01\%                        \\
\hline
  XSec Tot (corr)                  & 4.00\%                        \\
  XSec Tot                       & 4.08\%                        \\
\hline
\hline
  Flux+XSec (ND280 constrained)                      & 2.79\%                        \\
  Flux+XSec  (All)                    & 2.90\%                        \\
\hline
  Flux+XSec+SKDet+FSI+SI         & 5.03\%                        \\
  Flux+XSec+SKDet+FSI+SI (pre-fit)         & 12.0\%                        \\
\hline
\hline
  Oscillations                   & 0.03\%                        \\
\hline
\hline
  All                            & 5.06\%                        \\
  All (pre-fit) 		& 12.0\%				\\
\hline

\end{tabular}
\caption{
 Effect of 1$\sigma$ variations on the total number of events 1-ring muon-like events for neutrino mode
 given by all the systematic uncertainties, obtained by performing 10k toys MC.
The RMS / Mean of the distribution is assumed to be the $1\sigma$ uncertainty.
The systematic parameter values measured with the ND280 data 
shown in 
\cref{tab:banff_uncertainties_flux_fhc},
~\ref{tab:banff_uncertainties_flux_rhc}
~\ref{tab:banff_uncertainties_xsec}
and
~\ref{tab:banff_uncertainties_skdetfsi}
and the oscillation parameters of the Asimov data set A shown in \cref{tab:osc_params_joint_asimov}
are used.
Since the parameters $\delta_{CP}$, $\theta_{23}$ and $\Delta m^2_{32}$ ($\Delta m^2_{13}$) and mass hierarchy
will be measured,
only the uncertainty on the parameters 
$\sin^2 \theta_{13}$,
$\sin^2 \theta_{12}$
and 
$\Delta m^2_{21}$ 
are taken into account.
Pre-BANFF fit systematic uncertainties are used were stated.
The changes due to the Run 1-7c flux and covariance matrix updates are applied.
}                                                                                                                                                                                          
\label{tab:nsk_systematic_table_summary_joint_asimov_A_fhc_mulike_Run17c}
\end{table}




%%%%% fhc e-like


\begin{table}[htb]
\centering
\begin{tabular}{| c | c | c |}
\hline
  {\bf Source of uncertainty}    & \multicolumn{1}{c|}{{\bf ${\delta}N_{SK}$/$N_{SK}$}}\\
                                 & \multicolumn{1}{c|}{(Asimov A)}\\

\hline
\hline  
SKDet+FSI+SI                   & 3.46\%                        \\
SKDet only                   & 2.39\%                        \\
FSI+SI only                   & 2.50\%                        \\
\hline
\hline
  Flux                           & 3.64\%                        \\
\hline
\hline
  2p-2h (corr)                     & 3.87\%                        \\
  2p-2h bar (corr)                  & 0.05\%                        \\
  NC other (uncorr)              & 0.16\%                        \\
  NC 1gamma (uncorr)             & 1.44\%                        \\
  XSec nue/numu (uncorr)                 & 2.65\%                        \\
\hline
  XSec Tot (corr)                  & 4.13\%                        \\
  XSec Tot                       & 5.12\%                        \\
\hline
\hline
  Flux+XSec (ND280 constrained)                      & 2.88\%                        \\
  Flux+XSec  (All)                     & 4.17\%                        \\
\hline
  Flux+XSec+SKDet+FSI+SI         & 5.41\%                        \\
  Flux+XSec+SKDet+FSI+SI (pre-fit)         & 11.9\%                        \\
\hline
\hline
  Oscillations                   & 4.13\%                        \\
\hline
\hline
  All                            & 6.86\%                        \\
  All (pre-fit) 		& 12.8\%				\\
\hline
\hline

\end{tabular}
\caption{
 Effect of 1$\sigma$ variations on the total number of events 1-ring electron-like events for neutrino mode
 given by all the systematic uncertainties, obtained by performing 10k toys MC.
The RMS / Mean of the distribution is assumed to be the $1\sigma$ uncertainty.
The systematic parameter values measured with the ND280 data 
shown in 
\cref{tab:banff_uncertainties_flux_fhc},
~\ref{tab:banff_uncertainties_flux_rhc}
~\ref{tab:banff_uncertainties_xsec}
and
~\ref{tab:banff_uncertainties_skdetfsi}
and the oscillation parameters of the Asimov data set A shown in \cref{tab:osc_params_joint_asimov}
are used.
Since the parameters $\delta_{CP}$, $\theta_{23}$ and $\Delta m^2_{32}$ ($\Delta m^2_{13}$) and mass hierarchy
will be measured,
only the uncertainty on the parameters 
$\sin^2 \theta_{13}$,
$\sin^2 \theta_{12}$
and 
$\Delta m^2_{21}$ 
are taken into account.
Pre-BANFF fit systematic uncertainties are used were stated.
The changes due to the Run 1-7c flux and covariance matrix updates are applied.
}                                                                                                                                                                                          
\label{tab:nsk_systematic_table_summary_joint_asimov_A_fhc_elike_Run17c}
\end{table}



%%%%%% rhc mu-like


\begin{table}[htb]
\centering
\begin{tabular}{| c | c | c |}
\hline
  {\bf Source of uncertainty}    & \multicolumn{1}{c|}{{\bf ${\delta}N_{SK}$/$N_{SK}$}}\\
                                 & \multicolumn{1}{c|}{(Asimov A)}\\

\hline
\hline  
SKDet+FSI+SI                   & 3.90\%                        \\
SKDet only                   & 3.31 \%                        \\
FSI+SI only                   & 2.06 \%                        \\
\hline
  Flux                           & 3.77\%                        \\
\hline
\hline
  2p-2h (corr)                     & 2.96\%                        \\
  2p-2h bar (corr)                  & 1.81\%                        \\
  NC other (uncorr)              & 0.75\%                        \\
  NC 1gamma (uncorr)             & 0.00\%                        \\
  XSec nue/numu (uncorr)                 & 0.00\%                        \\
\hline
  XSec Tot (corr)                  & 4.13\%                        \\
  XSec Tot                       & 4.19\%                        \\
\hline
\hline
  Flux+XSec (ND280 constrained)                      & 3.26\%                        \\
  Flux+XSec  (All)                    & 3.35\%                        \\
\hline
  Flux+XSec+SKDet+FSI+SI         & 5.22\%                        \\
  Flux+XSec+SKDet+FSI+SI (pre-fit)         & 12.5\%                        \\
\hline
\hline
  Oscillations                   & 0.03\%                        \\
\hline
\hline
  All                            & 5.19\%                        \\
  All (pre-fit) 		& 12.5\%				\\
\hline
\hline
\end{tabular} 
\caption{
 Effect of 1$\sigma$ variations on the total number of events 1-ring muon-like events for antineutrino mode
 given by all the systematic uncertainties, obtained by performing 10k toys MC.
The RMS / Mean of the distribution is assumed to be the $1\sigma$ uncertainty.
The systematic parameter values measured with the ND280 data 
shown in 
\cref{tab:banff_uncertainties_flux_fhc},
~\ref{tab:banff_uncertainties_flux_rhc}
~\ref{tab:banff_uncertainties_xsec}
and
~\ref{tab:banff_uncertainties_skdetfsi}
and the oscillation parameters of the Asimov data set A shown in \cref{tab:osc_params_joint_asimov}
are used.
Since the parameters $\delta_{CP}$, $\theta_{23}$ and $\Delta m^2_{32}$ ($\Delta m^2_{13}$) and mass hierarchy
will be measured,
only the uncertainty on the parameters 
$\sin^2 \theta_{13}$,
$\sin^2 \theta_{12}$
and 
$\Delta m^2_{21}$ 
are taken into account.
Pre-BANFF fit systematic uncertainties are used were stated.
The changes due to the Run 1-7c flux and covariance matrix updates are applied.
}                                                                                                                                                                                          
\label{tab:nsk_systematic_table_summary_joint_asimov_A_rhc_mulike_Run17c}
\end{table}



%%%%%% rhc e-like

\begin{table}[htb]
\centering
\begin{tabular}{| c | c | c |}
\hline
  {\bf Source of uncertainty}    & \multicolumn{1}{c|}{{\bf ${\delta}N_{SK}$/$N_{SK}$}}\\
                                 & \multicolumn{1}{c|}{(Asimov A)}\\

\hline
\hline  
SKDet+FSI+SI                   & 3.95\%                        \\
SKDet only                   & 3.09\%                        \\
FSI+SI only                   & 2.46\%                        \\
\hline
  Flux                           & 3.77\%                        \\
\hline
\hline
  2p-2h (corr)                     & 2.97\%                        \\
  2p-2h bar (corr)                  & 2.36\%                        \\
  NC other (uncorr)              & 0.33\%                        \\
  NC 1gamma (uncorr)             & 2.95\%                        \\
  XSec nue/numu (uncorr)                 & 1.50\%                        \\
\hline
  XSec Tot (corr)                  & 4.32\%                        \\
  XSec Tot                       & 5.45\%                        \\
\hline
\hline
  Flux+XSec (ND280 constrained)                      & 3.22\%                        \\
  Flux+XSec                      & 4.63\%                        \\
\hline
  Flux+XSec+SKDet+FSI+SI         & 6.19\%                        \\
  Flux+XSec+SKDet+FSI+SI (pre-fit)         & 13.5\%                        \\
\hline
\hline
  Oscillations                   & 4.06\%                        \\
\hline
\hline
  All                            & 7.39\%                        \\
      All (pre-fit) 		& 14.1\%				\\
\hline
\hline
\end{tabular}
\caption{
 Effect of 1$\sigma$ variations on the total number of events 1-ring electron-like events for antineutrino mode
 given by all the systematic uncertainties, obtained by performing 10k toys MC.
The RMS / Mean of the distribution is assumed to be the $1\sigma$ uncertainty.
The systematic parameter values measured with the ND280 data 
shown in 
\cref{tab:banff_uncertainties_flux_fhc},
~\ref{tab:banff_uncertainties_flux_rhc}
~\ref{tab:banff_uncertainties_xsec}
and
~\ref{tab:banff_uncertainties_skdetfsi}
and the oscillation parameters of the Asimov data set A shown in \cref{tab:osc_params_joint_asimov}
are used.
Since the parameters $\delta_{CP}$, $\theta_{23}$ and $\Delta m^2_{32}$ ($\Delta m^2_{13}$) and mass hierarchy
will be measured,
only the uncertainty on the parameters 
$\sin^2 \theta_{13}$,
$\sin^2 \theta_{12}$
and 
$\Delta m^2_{21}$ 
are taken into account.
Pre-BANFF fit systematic uncertainties are used were stated.
The changes due to the Run 1-7c flux and covariance matrix updates are applied.
}                                                                                                                                                                                          
\label{tab:nsk_systematic_table_summary_joint_asimov_A_rhc_elike_Run17c}
\end{table}



 


\begin{table}
\center
\begin{tabular}{| c || c | c || c | c |}
\hline
Source of uncertainty 
&$\mu$-like ${\delta} \Big( \frac{\# {\nu} \text{-mode}}{\# {\bar{\nu}} \text{-mode} } \Big) $ / $\Big \langle \frac{\# {\nu} \text{-mode}}{\# {\bar{\nu}} \text{-mode} } \Big \rangle$ 
&$e$-like ${\delta} \Big( \frac{\# {\nu} \text{-mode}}{\# {\bar{\nu}} \text{-mode} } \Big) $ / $\Big \langle \frac{\# {\nu} \text{-mode}}{\# {\bar{\nu}} \text{-mode} } \Big \rangle$ \\


\hline
\hline
SKDet & 0.07\% & 1.6\% \\ \hline
FSI+SI & 2.6\% & 3.6\% \\ \hline
SKDet+FSI+SI & 2.6\% & 3.8\% \\ \hline
\hline
Flux & 1.8\% & 1.8\% \\ \hline
XSec NC other (uncorr) & 0.0\% & 0.2\% \\ \hline
XSec NC $1\gamma$ (uncorr) & 0.0\% & 1.5\% \\ \hline
XSec $\nu_e$ / $\nu_{\mu}$ (uncorr) & 0.0\% & 3.1\% \\ \hline
XSec Tot & 1.7\% & 4.3\% \\ \hline
 \hline
Flux+XSec & 1.9\% & 4.1\% \\ \hline
 \hline
Flux+XSec+SKDet+FSI+SI & 3.3\% & 5.7\% \\ \hline
 \hline
Oscillations & 0.0\% & 0.8\% \\ \hline
 \hline
All & 3.2\% & 5.8\% \\ \hline

%\hline
%\hline
%SKDet 							& XXX\%  	& XXX\% \\ \hline
%FSI+SI 	 						& XXX\% 	& XXX\% \\ \hline
%SKDet+FSI+SI 						& XXX\% 	& XXX\% \\ \hline
%\hline
%Flux 								& XXX\% 	& XXX\% \\ \hline
%XSec NC other (uncorr)  				& XXX\% 	& XXX\% \\ \hline
%XSec NC $1\gamma$ (uncorr) 			& XXX\%  	& XXX\% \\ \hline
%XSec $\nu_e$ / $\nu_{\mu}$ (uncorr)  	& XXX\%  	& XXX\% \\ \hline
%XSec Tot (All)						& XXX\%  	& XXX\% \\ \hline
% \hline
%Flux+XSec 						& XXX\%  	& XXX\% \\ \hline
% \hline
%Flux+XSec+SDet+FSI+SI 				& XXX\%  	& XXX\% \\ \hline
% \hline
%Oscillations 						& XXX\% 	& XXX\% \\ \hline
% \hline
%All 								& XXX\%  	& XXX\% \\ \hline

\end{tabular}
\caption{
 Effect of 1$\sigma$ variations of the several source of systematic uncertainties 
 on the ratio between the
 total number of events in neutrino and antineutrino modes
 for 1-ring electron- and muon-like samples
 obtained by performing 10k toys MC.
 The systematic parameter values measured with the ND280 data 
shown in 
\cref{tab:banff_uncertainties_flux_fhc},
~\ref{tab:banff_uncertainties_flux_rhc}
~\ref{tab:banff_uncertainties_xsec}
and
~\ref{tab:banff_uncertainties_skdetfsi}
and the oscillation parameters of the Asimov data set A shown in \cref{tab:osc_params_joint_asimov}
are used.
Since the parameters $\delta_{CP}$, $\theta_{23}$ and $\Delta m^2_{32}$ ($\Delta m^2_{13}$) and mass hierarchy
will be measured,
only the uncertainty on the parameters 
$\sin^2 \theta_{13}$,
$\sin^2 \theta_{12}$
and 
$\Delta m^2_{21}$ 
are taken into account.
The changes due to the Run 1-7c flux and covariance matrix updates are applied.
}             
\label{tab:ratio_systematic_table_summary_asimov_A_Run17c}                                                                                                                                                                             
\end{table}
